\begin{pastexam}[label={px:ncku-112-ee-ds-08}]{成大電機所 112 學年度 第 8 題}
（問答）

(15\%) The following puzzle is called a word ladder: In a word ladder puzzle, you must make the
change occur gradually by changing one letter at a time. At each step you must transform one word
into another word, you are not allowed to transform a word into a non-word. The following sequence
of words shows one possible solution to the problem ``Transform the word ``FOOL'' into the word
``SAGE''.''

\hspace*{1em}Example: FOOL -> POOL -> POLL -> POLE-> PALE -> SALE -> SAGE

Here, we are interested in figuring out the smallest number of transformations needed to turn the
starting word into the ending word.

(a) Construct an example graph by using these words: FOOL, POOL, POLL, POLE, PALE, SALE, SAGE, COOL,
FOUL, FOIL, FAIL, FALL, PALL, POLE, POPE, PAGE. Represent the relationships between the words as a
word ladder graph. (Hint: Think about the path where we get from the word ``FOOL'' to the word
``SAGE'').

(b) Discuss which search algorithm you can use to find an efficient path from the starting word to
the ending word, why it is more efficient compared with other. Discuss the complexity.

(c) There is one way to improve efficiency instead of using a graph to show the relationship.
Suppose that we have a huge number of \textbf{buckets}, each of them with a four-letter word on the
outside, except that one of the letters in the label has been replaced by an underscore. The labels
on the buckets are the keys in our dictionary. The value stored for that key is a list of words.
\textbf{Show the data structure and list the buckets using the following words} (don't list the
buckets containing only one item): FOOL, POOL, POLL, POLE, PALE, SALE, SAGE, COOL, FOUL, FOIL, FAIL,
FALL, PALL, POLE, POPE, PAGE
\end{pastexam}

\begin{pastsoln}
本題和成大智慧資安學程同年的第~8~題相同（電機所、電通所、智慧資安學程三個系所共用同一份考卷），完整的詳解見歷屆試題~\ref{px:ncku-112-isec-ds-08}，這裡只列答案和重點。

原題的單字表把 POLE 列了兩次，同一個字只算一個頂點，共 $15$ 個字。

答案：
\begin{itemize}
\item (a) 每個字一個頂點，只差一個字母的兩個字連一條無向邊，共 $19$ 條邊（鄰接表和圖見智慧資安學程的
那一題）。
\item (b) 用 BFS（資結書第~21.3~節）。 每一步代價相同，BFS 依距離一層一層擴展，第一次碰到終點時的
層數就是最少變換次數。 FOOL 到 SAGE 最少 $6$ 步，例如 FOOL $\to$ POOL $\to$ POLL $\to$ PALL
$\to$ PALE $\to$ PAGE $\to$ SAGE（題目的範例也是 $6$ 步）。 DFS 第一次找到的路不一定最短；
Dijkstra 可以但要優先佇列，多一個 $\log$ 因子。 圖建好後 BFS 是 $O(|V| + |E|)$。
\item (c) 用字典（雜湊表）：鍵是把一個位置換成底線的標籤，值是符合的字的串列；每個字放進 $4$ 個桶子。
兩個字以上的桶子有 $15$ 個：\texttt{\_OOL}：FOOL, POOL, COOL；\texttt{FO\_L}：FOOL, FOUL, FOIL；
\texttt{PO\_L}：POOL, POLL；\texttt{P\_LL}：POLL, PALL；\texttt{POL\_}：POLL, POLE；
\texttt{P\_LE}：POLE, PALE；\texttt{PO\_E}：POLE, POPE；\texttt{PA\_E}：PALE, PAGE；
\texttt{PAL\_}：PALE, PALL；\texttt{SA\_E}：SALE, SAGE；\texttt{\_AGE}：SAGE, PAGE；
\texttt{F\_IL}：FOIL, FAIL；\texttt{FA\_L}：FAIL, FALL；\texttt{\_ALL}：FALL, PALL；\texttt{\_ALE}：PALE, SALE。
同一個桶子裡的字兩兩相鄰，建圖只要把每個字放進 $4$ 個桶子，不必兩兩比較所有的字（資結書第~12~章）。
\end{itemize}
\end{pastsoln}
